\documentclass[10pt,a4paper]{amsart}
\usepackage[margin=19mm]{geometry}
\usepackage{amsmath,amssymb,mathtools}
\usepackage[scr=rsfs,cal=euler]{mathalfa}
\usepackage{xfrac}
\usepackage[unicode,hidelinks,bookmarks=false]{hyperref}
\newcommand{\ie}{i.e.\ }
\newcommand{\cf}{cf.\ }
\newcommand{\resp}{respectively\ }
\newcommand{\Subsection}{\subsection}
\providecommand{\nobreakdash}{\nobreak\hbox{-}}
\setlength{\emergencystretch}{2em}
\multlinegap0pt
\usepackage{amssymb}
\usepackage{mathtools}
\usepackage{xcolor}
\usepackage[normalem]{ulem}
\newtheorem{lemma}{Lemma}
\newtheorem{remark}{Remark}
\newtheorem{proposition}{Proposition}
\usepackage{cleveref}
\crefname{lemma}{Lemma}{Lemmas}
\crefname{remark}{Remark}{Remarks}
\crefname{proposition}{Proposition}{Propositions}
\newcommand{\Q}{\mathbb{Q}}
\newcommand{\Z}{\mathbb{Z}}
\title{Explicit $p$-adic Hodge theory for elliptic curves and non-split Cartan images}
\author{Matthew Bisatt, Lorenzo Furio,, Davide Lombardo}
\date{Source: arXiv:2603.04021v2 (CC BY 4.0)}
\begin{document}
\maketitle

\noindent\emph{Source and licence.} Explicit $p$-adic Hodge theory for elliptic curves and non-split Cartan images. Version arXiv:2603.04021v2 (CC BY 4.0), distributed by arXiv under the Creative Commons Attribution 4.0 International licence, http://creativecommons.org/licenses/by/4.0/, as recorded in the arXiv licence metadata for this exact version and shown on the abstract page https://arxiv.org/abs/2603.04021v2. Full licence notice: \texttt{licenses/arxiv-2603.04021-CC-BY-4.0.txt}.\par\smallskip

\noindent\textit{Editorial scope.} Counted unit: the article's proposition congruence (source0.tex lines 1590--1596). The article's complete original proof of it is delivered with it (source0.tex lines 1598--1636, one \texttt{proof} environment of 39 lines). No mathematics is written, repaired or re-derived: the statement and the proof are the article's own lines, in the article's own notation, and no retained line is substituted or re-flowed.  Retained uncounted as the source-local context the proof needs: lines 1392--1397; lines 1442--1447; lines 1482--1487; lines 1508--1515; lines 1585--1587.  Nothing was omitted from any retained span, so the package carries no omission record.  No retained line contains a citation, so the wrapper carries no bibliography and no bibliography entry.  No retained line contains a figure environment, an image-inclusion command or drawing code, so no image is delivered and the package declares no asset dependency.  Editorial formatting only, in addition: an AMS \texttt{amsart} preamble that restates the article's own declarations and macros used by the retained lines, namely the environments \texttt{lemma}, \texttt{remark}, \texttt{proposition} on separate counters, together with the standard packages those lines need.  The retained environments print in this document as Lemma 1, Remark 1, Proposition 1 in order of appearance, whereas the article prints the same items with the numbering of its own full document.  The complete native source of the article as distributed in the arXiv e-print for this exact version is bundled unchanged as \texttt{sources/195712/source0.tex}.
\begin{lemma}\label{lemma: freshmans dream}
Let $n\geq 1$ be an integer and $a,b \in \overline{\Q}_p$ with $v(a) \geq 0, v(b) \geq 0$. Then
\[
(a+b)^{p^n} \equiv a^{p^n}+b^{p^n} \bmod p^{ 1 + p^n \min\{v(a), v(b)\}}.
\]
\end{lemma}

\begin{lemma}\label{lemma: congruences powers of gamma}
We have $\gamma^{p^{2k-2}(p^2-1)} \equiv p(1 + \alpha^{-1} \pi_e^2 \gamma^{p^{2k-2}(p-1)}) \mod p^{2-\frac{1}{p^2}}$. Moreover, there exists an $e$-th root of unity $\xi$ such that
\[
\gamma^{p^{2k-2} \cdot \frac{2(p^2-1)}{e}} \equiv \xi \pi_e^2 \left(1 + \frac{2}{e}\alpha^{-1}\pi_e^2 \gamma^{p^{2k-2}(p-1)} \right) \quad \bmod{p^{1+ \frac{2}{p+1}}}.
\]
\end{lemma}

\begin{lemma}\label{lemma: epsilon1}
Suppose $p > 3$. Let $\varepsilon_1=\beta-\zeta\gamma$ and $\alpha^{-1} = u p^{k-2}$ with $u \in \Z_p^\times$. Then $\varepsilon_1$ satisfies
\begin{equation}\label{eq: lemma epsilon1}
     (-1)^k\pi_e^2u(\zeta-\zeta^p)\gamma^{p^{2k-1}} - \varepsilon_1^{p^2} + p\varepsilon_1 \equiv 0 \quad \bmod{p^{1+\frac{2}{ep}+\frac{1}{p^2-1}}}. 
\end{equation}
\end{lemma}

\begin{lemma}\label{lemma: u1 and u2}
Let $u_1=\xi^{-1}\alpha^{-1}(\zeta-\zeta^p)$, $u_2=-\frac{2}{e} \,\xi^{-2}\alpha^{-2}(\zeta-\zeta^p)$, $h_1=\frac{2(p^2-1)}{e}+p$ and $h_2=\frac{4(p^2-1)}{e}+2p-1$. 
Then
\begin{enumerate}
    \item $(u_1\gamma^{h_1})^{p^2} \equiv \alpha^{-1}\pi_e^2(\zeta-\zeta^p)\gamma^{p^3}\left(1+\tfrac{2}{e}\alpha^{-1}\pi_e^2\gamma^{p^3-p^2}\right) \quad \bmod{p^{1+\frac{1}{p^2-1}}}.$
    \item $(u_2\gamma^{h_2})^{p^2} \equiv -\frac{2}{e}\alpha^{-2}\pi_e^4(\zeta-\zeta^p)\gamma^{2p^3-p^2} \quad \bmod{p^{1+\frac{1}{p^2-1}}}.$
\end{enumerate}
\end{lemma}

\begin{remark}\label{rmk: h_i inequalities}
    Note that $h_i=\frac{1}{p^{2i-2}}(h_1-p^{2k-2}) + p^{2k-2}$ for $i \geq 1$. As $h_1 < p^{2k-2}$, it is immediate that the sequence $h_i$ is strictly increasing and bounded by $p^{2k-2}$. Moreover, $h_1 \in \Z$ and $v(h_i)=2(k-i-1)$ so $h_1,\ldots,h_{k-1}$ are integers, as is $p^2h_k$. In particular, this implies $\varepsilon_i$ are well-defined elements of the splitting field of $g_k$ for all $1 \leq i \leq k$.
\end{remark}

\begin{proposition}\label{prop: recursive congruence}
    Suppose $p>3$. With the notation above, we have
    \[
    v(\varepsilon_i) = h_iv(\gamma) \quad \text{for all } 1 \leq i \leq k.
    \]
    In particular $v(\varepsilon_k)$ has denominator divisible by $p^{2k}$.
\end{proposition}

\begin{proof}
    We first claim that the congruence
    \begin{equation}\label{eq: eps congruence}
    u_1 \gamma^{p^2h_i} - \varepsilon_i^{p^2} + p\varepsilon_i \equiv 0 \mod p^{1+\frac{2}{ep} + \frac{1}{p^2-1}}
    \end{equation}
    holds for every $1 \leq i \leq k$ (note the exponents $p^2h_i$ are integers by \Cref{rmk: h_i inequalities}). We defer the proof of this momentarily, and instead first use it to deduce the valuation of $\varepsilon_i$. 
    By  \Cref{rmk: h_i inequalities}, we have $h_i < p^{2k-2}$ and hence $v(\gamma^{p^2h_i})<1+\frac{1}{p^2-1}$. If $v(\varepsilon_i) \geq \frac{1}{p^2-1}$, then
    \[
    v(\varepsilon_i^{p^2}) \geq v(p\varepsilon_i) \geq 1 + \frac{1}{p^2-1} > v(u_1\gamma^{p^2h_i})
    \]
    which cannot satisfy the congruence \eqref{eq: eps congruence} (noting that $1 + \frac{1}{p^2-1}< 1 + \frac{2}{ep}+\frac{1}{p^2-1}$) and hence $v(\varepsilon_i) < \frac{1}{p^2-1}$. Therefore $v(\varepsilon_i^{p^2})< v(p\varepsilon_i)<1+\frac{1}{p^2-1}$ so to obtain the cancellation we must have $v(\varepsilon_i^{p^2})=v(u_1\gamma^{p^2h_i})$, equivalently $v(\varepsilon_i)=h_iv(\gamma)$.
    In particular, if the congruence is true for $\varepsilon_k$, then since $v(h_k)=-2$ by \Cref{rmk: h_i inequalities}, we have
    \[
    v(v(\varepsilon_k))=v(h_kv(\gamma))=v(h_k)+v(v(\gamma))=-2 -(2k-2) = -2k
    \] 
    so the denominator of $v(\varepsilon_k)$ is divisible by $p^{2k}$ as required.
    
    It remains to prove the claimed congruence \eqref{eq: eps congruence}, which we do by induction on $i \leq k$. Suppose first that $i = 1$. By \Cref{lemma: epsilon1} we have
    \[
    (-1)^k\pi_e^2u(\zeta-\zeta^p)\gamma^{p^{2k-1}} - \varepsilon_1^{p^2} + p\varepsilon_1 \equiv 0 \mod{p^{1+\frac{2}{ep}+\frac{1}{p^2-1}}},
    \]
    hence it suffices to show that 
    \[
    u_1 \gamma^{p^2h_1} := (-1)^k\xi^{-1}u(\zeta-\zeta^p) \gamma^{p^{2k-2} \cdot\frac{2(p^2-1)}{e} + p^{2k-1}} \equiv (-1)^k\pi_e^2u(\zeta-\zeta^p)\gamma^{p^{2k-1}} \mod p^{1+\frac{2}{ep}+\frac{1}{p^2-1}}.
    \]
    The base case now follows from \Cref{lemma: congruences powers of gamma}, noting that $v(\alpha^{-1}) \ge 1$ (as $k>2$), and that $1+ \frac{2}{ep} + \frac{1}{p^2-1} < 1 + \frac{2}{p+1}$.

    We now suppose the congruence holds for $\varepsilon_i$ (so in particular $v(\varepsilon_i)=h_iv(\gamma)$) and prove the induction step: namely that the congruence also holds for $\varepsilon_{i+1}$ whenever $i \leq k-1$. 
    Recall that $\varepsilon_{i+1}=\varepsilon_{i}-u_1\gamma^{h_i}$ by definition. Using that $v(\varepsilon_{i})=h_iv(\gamma)$ and $h_i \geq h_1$ (\Cref{rmk: h_i inequalities}) we have
    \[
    p^2v(\varepsilon_{i+1}) \geq p^2v(\gamma^{h_i}) \geq p^2v(\gamma^{h_1}) \geq \frac{2}{e}+\frac{1}{p^2-1}.
    \]
    The freshman's dream (\Cref{lemma: freshmans dream}) therefore gives us that $\varepsilon_{i}^{p^2} \equiv \varepsilon_{i+1}^{p^2} + u_1^{p^2}\gamma^{p^2h_i} \mod p^{1+\frac{2}{ep}+\frac{1}{p^2-1}}$. Moreover, analogously to the first line of the proof of \Cref{lemma: u1 and u2}, $u_1^{p^2}\equiv u_1 \bmod p$, and hence $u_1^{p^2} \gamma^{p^2 h_i} \equiv u_1 \gamma^{p^2h_i} \mod p^{1+\frac{2}{ep}+\frac{1}{p^2-1}}$ since $v(\gamma^{p^2h_i}) \ge\frac{2}{e}+\frac{1}{p^2-1} > \frac{2}{ep} + \frac{1}{p^2-1}$ by above.
    Rewriting the congruence for $\varepsilon_i$ now yields
    \begin{equation*}
        0 \equiv u_1\gamma^{p^2h_i} - (\varepsilon_{i+1} + u_1\gamma^{h_i})^{p^2} + p\varepsilon_{i+1} + pu_1\gamma^{h_i} \equiv p u_1 \gamma^{h_i} - \varepsilon_{i+1}^{p^2} + p\varepsilon_{i+1} \mod p^{1+\frac{2}{ep}+\frac{1}{p^2-1}}.
    \end{equation*}
    To obtain the claimed congruence for $\varepsilon_{i+1}$, it remains to show that $pu_1 \gamma^{h_i} \equiv u_1\gamma^{p^2h_{i+1}} \mod p^{1+\frac{2}{ep}+\frac{1}{p^2-1}}.$ Since $v(\alpha^{-1}) \geq 1$, we have $p \equiv \gamma^{p^{2k-2}(p^2-1)} \mod p^{2 - \frac{1}{p^2}}$ by \Cref{lemma: congruences powers of gamma}, so substituting this in and using the recurrence relation for $h_i$ proves that $pu_1 \gamma^{h_i} \equiv u_1\gamma^{p^2h_{i+1}} \mod p^{1+\frac{2}{ep}+\frac{1}{p^2-1}}$ and hence also the claimed congruence \eqref{eq: eps congruence} as required.
\end{proof}

\end{document}
